\chapter*{Conclusion}
\addcontentsline{toc}{chapter}{Conclusion}
\markboth{Conclusion}{}

Ce rapport a appliqué, pour la première fois à notre connaissance, l'outil de modélisation financière d'une mine d'or diffusé par NRGI --- adaptation simplifiée de la méthodologie FARI du FMI --- au cas du secteur aurifère sénégalais, en s'appuyant sur les données techniques et économiques de l'étude de préfaisabilité du projet Sabodala-Massawa publiée en 2020. Deux opérations ont été nécessaires~: remplacer le profil opérationnel générique de l'outil par le profil réel du projet (production, dépenses d'investissement, coûts d'exploitation, prix de l'or) sur la période 2020-2036, et compléter le tableau comparatif des régimes fiscaux --- qui comportait dix régimes pré-paramétrés et une colonne vide --- par un onzième régime reproduisant les paramètres effectivement applicables à Sabodala Gold Operations (redevance de 5\,\% sur la valeur des expéditions, impôt sur les sociétés de 25\,\%, participation gratuite de l'État de 10\,\%, exonération de droits de douane).

Les résultats indiquent que, pour ce profil de projet, le régime fiscal sénégalais génère un \TEMI{} de 44,3\,\%, sensiblement inférieur à celui du code minier congolais amendé (57,6\,\%) auquel l'outil était initialement consacré, et proche du bas de la distribution des onze régimes comparés --- seuls la Tanzanie et le Pérou présentant un \TEMI{} inférieur pour ce même profil. La structure des recettes publiques sénégalaises est dominée par l'impôt sur les bénéfices (49\,\%), suivi de la redevance (22\,\%), de la participation gratuite de l'État (16\,\%) et de la retenue à la source sur les dividendes (14\,\%)~; l'absence de droits de douane, liée au statut d'entreprise franche d'exportation de SGO, constitue l'écart le plus significatif avec les régimes de comparaison.

Un second résultat, de nature davantage méthodologique, concerne les indicateurs de rentabilité du point de vue de l'investisseur~: le \TRI{} avant et après impôt, ainsi que les périodes de récupération, ne sont pas interprétables pour Sabodala-Massawa, dans la mesure où le flux de trésorerie avant impôt du projet est positif dès la première année du modèle --- caractéristique d'un projet d'extension d'une mine déjà en production, et non d'un projet \emph{greenfield} pour lequel ces indicateurs ont été conçus. La \VAN{} et le \TEMI, en revanche, demeurent pleinement valides et constituent, à ce titre, les indicateurs les plus pertinents pour ce type de configuration.

Ces résultats appellent trois prolongements. D'une part, l'élargissement de l'exercice aux autres titres miniers en exploitation au Sénégal --- notamment dans le secteur du zircon et des phosphates --- permettrait de disposer d'une vue d'ensemble de la performance fiscale du secteur extractif, et non d'un seul projet. D'autre part, une analyse de sensibilité au prix de l'or, qui n'a pas été conduite ici par souci de cohérence avec le cas central retenu par le PFS, donnerait une mesure de la stabilité du \TEMI{} sénégalais sur le cycle des prix, complémentaire de la photographie statique présentée au chapitre~\ref{chap:resultats}. Enfin, la quantification --- même approximative --- des flux contractuels qui échappent au régime fiscal général, tels que le contrat d'approvisionnement en or conclu avec Franco-Nevada dans le cadre de l'acquisition de Massawa, permettrait de rapprocher le \TEMI{} « modélisé » du partage effectif de la valeur entre l'État sénégalais et les actionnaires du projet.

Au-delà de ces résultats propres à Sabodala-Massawa, ce rapport illustre la possibilité, pour une administration disposant de ressources d'analyse limitées, de s'appuyer sur des outils ouverts --- ici, le classeur NRGI --- pour produire des estimations chiffrées, reproductibles et actualisables de la performance fiscale d'un projet minier, sans remettre en cause l'intégrité méthodologique de ces outils. C'est dans cette perspective que les pistes formulées au chapitre~\ref{chap:recommandations} ont été proposées pour le travail futur du Comité National ITIE Sénégal.
